% Provenance (never printed): SNT "Le Web" (URL, HTTP, browser, client-server, static/dynamic pages, HTTPS) and NSI
% premiere "Interactions entre l'homme et la machine sur le Web" (events, forms, GET/POST, client and server
% roles, sessions). See sources/book2/BRIEFS.md.
\chapter{The Web: Clients, Servers and HTTP}\label{ch:g10:http-clients-servers}

Sara types the address of the town library's site and presses Enter. Less than a second later the page is there,
with its title, a logo, a search box and two small pictures. In that second her browser and a machine she has never
seen exchanged about a dozen messages, each a few lines of plain text, written in an agreed format.

This chapter reads those messages. It runs a small web server on the reader's own machine, at the address
\code{127.0.0.1}, which always means ``this computer'', and sends it \omterm{def:g10:http-clients-servers:http}{requests} from Python, so that every line
that crosses between the two can be printed.

\begin{recall}
Data crosses the internet in packets; every machine has an address, and the name system turns names into addresses.
A client asks and a server answers (Book 1, grade 8). A browser shows web pages, and a link takes you to another
page (Book 1, grade 5). A page is written in HTML, and its look in CSS; the padlock next to an address means the
connection is encrypted (Book 1, grade 9).
\end{recall}

\section{The parts of an address}\label{sec:g10:http-clients-servers:url}

% ledger: g10:uri-3986
\begin{definition}[URL]\label{def:g10:http-clients-servers:url}
A \emph{URL}\index{URL} (uniform resource locator) is the address of a resource on the web. It is made of a
\emph{scheme}\index{scheme} (\code{http} or \code{https}), a \emph{host}\index{host} (the name or address of the
server), an optional \omterm{def:g10:http-clients-servers:port-number}{port number}, a path naming the resource on that server, an optional \emph{query
string}\index{query string} after a \code{?}, made of \code{name=value} pairs separated by \code{&}, and an optional
\emph{fragment}\index{fragment} after a \texttt{\#}, a place inside the resource that the browser keeps for itself.
\end{definition}

\begin{omfigure}
\begin{tikzpicture}[font=\ttfamily\small, every node/.style={inner sep=0pt, outer sep=0pt}]
  \node[anchor=base west] (a) at (0,0) {http};
  \node[anchor=base west] (b) at (a.base east) {://};
  \node[anchor=base west] (c) at (b.base east) {library.brenmoor.example};
  \node[anchor=base west] (d) at (c.base east) {:8080};
  \node[anchor=base west] (e) at (d.base east) {/books/search};
  \node[anchor=base west] (f) at (e.base east) {?title=forest\&genre=nature};
  \node[anchor=base west] (g) at (f.base east) {\#top};
  \begin{scope}[font=\footnotesize\rmfamily, every node/.style={inner sep=1pt, anchor=north}]
    \foreach \n/\lab/\col in {a/scheme/omDef, c/host/omProp, d/port/omThm, e/path/omMeth, f/query string/omDef,
        g/fragment/omProp}
      \draw[decorate, decoration={brace, amplitude=3pt, mirror}, \col, thick] ([yshift=-3pt]\n.south west) --
        ([yshift=-3pt]\n.south east) node[midway, yshift=-4pt, text=\col] {\lab};
  \end{scope}
\end{tikzpicture}

{\small The parts of an address. Without a \omterm{def:g10:http-clients-servers:port-number}{port number}, the \omterm{def:g10:http-clients-servers:url}{scheme} gives the usual one: $80$ for \code{http},
$443$ for \code{https}.}
\end{omfigure}

\omcode{code/grade-10/08-http-clients-servers/listings/parse_url.py}{1}{12}{Splitting an
address.}\label{lst:g10:http-clients-servers:parse-url}

\omoutput{code/grade-10/08-http-clients-servers/outputs/parse_url.txt}{Each part on its own; an address without a
\omterm{def:g10:http-clients-servers:port-number}{port number} has \code{None}.}\label{out:g10:http-clients-servers:parse-url}

The \omterm{def:g10:http-clients-servers:url}{host} names a machine, but one machine can run several servers at once: a web server, a mail server, a game
server. The \omterm{def:g10:http-clients-servers:port-number}{port number} tells them apart.

% ledger: g10:http-ports
\begin{definition}[Port number]\label{def:g10:http-clients-servers:port-number}
A \emph{port number}\index{port number} is a whole number from $0$ to $65{,}535$ that identifies one program waiting
for connections on a machine. A connection goes to a pair (address, port). Web servers usually listen on port $80$
for \code{http} and $443$ for \code{https}.
\end{definition}

\section{A request and its response}\label{sec:g10:http-clients-servers:http}

\begin{definition}[Protocol]\label{def:g10:http-clients-servers:protocol}
A \emph{protocol}\index{protocol} is a set of rules for an exchange of messages between programs: who speaks first,
the format of each message, and what each answer means.
\end{definition}

% ledger: g10:http-9112-message
\begin{definition}[HTTP, request, response]\label{def:g10:http-clients-servers:http}
\emph{HTTP}\index{HTTP} (hypertext transfer \omterm{def:g10:http-clients-servers:protocol}{protocol}) is the \omterm{def:g10:http-clients-servers:protocol}{protocol} of the web. The client sends a
\emph{request}\index{request}; the server sends back one \emph{response}\index{response}. Each message is text: a
first line, then \omterm{def:g10:http-clients-servers:http-header}{header} lines, then an empty line, then an optional body. Every line ends with the two characters
carriage return and line feed, written \texttt{\textbackslash r\textbackslash n} in Python.
\end{definition}

\begin{definition}[Header]\label{def:g10:http-clients-servers:http-header}
A \emph{header}\index{header (HTTP)} is a line \code{Name: value} of an \omterm{def:g10:http-clients-servers:http}{HTTP} message that gives information about
the message: the server's name (\code{Host}), the type of the body (\code{Content-Type}), its length in bytes
(\code{Content-Length}), a \omterm{def:g10:http-clients-servers:cookie}{cookie}, and so on.
\end{definition}

To see the messages themselves, the next listing opens a connection to the server and writes a \omterm{def:g10:http-clients-servers:http}{request} byte by byte,
as a browser would, then prints every line of the \omterm{def:g10:http-clients-servers:http}{response}.

\omcode{code/grade-10/08-http-clients-servers/listings/raw_exchange.py}{1}{15}{A \omterm{def:g10:http-clients-servers:http}{request} written by hand, and the
\omterm{def:g10:http-clients-servers:http}{response} line by line.}\label{lst:g10:http-clients-servers:raw-exchange}

\omoutput{code/grade-10/08-http-clients-servers/outputs/raw_exchange.txt}{The status line, three \omterm{def:g10:http-clients-servers:http-header}{headers}, the empty
line, the body. The last line is the server's own log, printed when it stops.}\label{out:g10:http-clients-servers:raw-exchange}

The helper \code{running}, in the chapter's folder, starts the server in a process of its own on the given port,
waits until it accepts connections, and stops it at the end of the \code{with} block, printing what the server wrote
meanwhile. The \code{Date} \omterm{def:g10:http-clients-servers:http-header}{header} is left out of the printout because it changes every second.

\begin{omfigure}
\begin{tikzpicture}[font=\footnotesize]
  \node[draw, fill=omDef!8, align=left, font=\ttfamily\footnotesize, inner sep=4pt] (q) {%
    GET /hello HTTP/1.1\\Host: 127.0.0.1\\Connection: close\\\ \\};
  \node[draw, fill=omProp!10, align=left, font=\ttfamily\footnotesize, inner sep=4pt, right=26mm of q] (r) {%
    HTTP/1.1 200 OK\\Server: BrenmoorSite/1.0\\Content-Type: text/plain; charset=utf-8\\Content-Length: 20\\\ \\
    Hello from Brenmoor};
  \node[above=1mm of q, font=\footnotesize\bfseries] {request};
  \node[above=1mm of r, font=\footnotesize\bfseries] {response};
  \draw[ptr] ([yshift=2mm]q.east) -- node[above] {client to server} ([yshift=2mm]q.east -| r.west);
  \draw[ptr, omProp] ([yshift=-4mm]q.east -| r.west) -- node[below] {server to client} ([yshift=-4mm]q.east);
  \node[left=1mm of q.north west, anchor=north east, align=right, text=omDef] {request line\\headers\\\\empty line};
\end{tikzpicture}

{\small One exchange. The \omterm{def:g10:http-clients-servers:http}{request} line names a \emph{\omterm{def:g10:http-clients-servers:methods}{method}}, a path and the version of \omterm{def:g10:http-clients-servers:http}{HTTP}; the status line gives
the version, a \omterm{def:g10:http-clients-servers:status}{status code} and its meaning in words. The body follows the empty line, and \code{Content-Length}
tells the client where it ends.}
\end{omfigure}

Here is the server. Its code is grouped in a \emph{class}, a notion of grade 12; for now, read it as a set of
functions that the \code{http.server} module calls: \code{do_GET} for each \omterm{def:g10:http-clients-servers:methods}{GET} \omterm{def:g10:http-clients-servers:http}{request}, \code{do_POST} for each \omterm{def:g10:http-clients-servers:methods}{POST}
\omterm{def:g10:http-clients-servers:http}{request}, and \code{self} as the way to reach the \omterm{def:g10:http-clients-servers:http}{request} and to answer it.

The file starts by importing \code{http.server} and \code{urllib.parse}, and by naming the folder of the site's
files and the type of each kind of file; then comes the class.

\omcode{code/grade-10/08-http-clients-servers/listings/library_site.py}{11}{44}{The library site: answering a \omterm{def:g10:http-clients-servers:methods}{GET}
\omterm{def:g10:http-clients-servers:http}{request}.}\label{lst:g10:http-clients-servers:site-get}

Python's \code{urllib.request} module is a client: it writes the \omterm{def:g10:http-clients-servers:http}{request}, reads the \omterm{def:g10:http-clients-servers:http}{response} and separates its
parts.

\omcode{code/grade-10/08-http-clients-servers/listings/fetch_once.py}{1}{13}{A client with
\code{urllib}.}\label{lst:g10:http-clients-servers:fetch-once}

\omoutput{code/grade-10/08-http-clients-servers/outputs/fetch_once.txt}{The same exchange, seen by a
client library.}\label{out:g10:http-clients-servers:fetch-once}

% ledger: g10:http-9110-classes
\begin{definition}[Status code]\label{def:g10:http-clients-servers:status}
The \emph{status code}\index{status code} of a \omterm{def:g10:http-clients-servers:http}{response} is a three-digit number whose first digit gives its class:
\code{1xx} information, \code{2xx} success (\code{200 OK}), \code{3xx} redirection, the resource is elsewhere
(\code{303 See Other}), \code{4xx} an error of the client (\code{404 Not Found}), \code{5xx} an error of the server.
\end{definition}

\omcode{code/grade-10/08-http-clients-servers/listings/status_codes.py}{1}{12}{Three paths, three
outcomes.}\label{lst:g10:http-clients-servers:status-codes}

\omoutput{code/grade-10/08-http-clients-servers/outputs/status_codes.txt}{The client follows the redirection on its
own: the server's log shows two \omterm{def:g10:http-clients-servers:http}{requests} for \code{/old-hello}.}\label{out:g10:http-clients-servers:status-codes}

A browser does the same, and does more: once the page has arrived, it reads its HTML and sends one more \omterm{def:g10:http-clients-servers:http}{request} for
each stylesheet and each image the page names.

\begin{omfigure}
\begin{tikzpicture}[y=-7mm, font=\footnotesize]
  \draw[thick] (0,0) node[above] {browser} -- (0,7.6);
  \draw[thick] (8,0) node[above] {server} -- (8,7.6);
  \foreach \y/\q/\r in {0.4/{GET /page.html}/{200, 425 bytes of HTML}, 2.6/{GET /style.css}/{200, 92 bytes},
      4.8/{GET /logo.svg}/{200, 204 bytes}} {
    \draw[ptr] (0,\y) -- node[pos=0.3, above, sloped] {\q} (8,\y+0.6);
    \draw[ptr, omProp] (8,\y+0.9) -- node[pos=0.3, below, sloped] {\r} (0,\y+1.5);
  }
  \node[align=center] at (4,7.2) {\dots\ and two more images};
  \node[anchor=west, align=left, text=black!60] at (8.3,1.3) {the browser reads the HTML\\and finds 4 resources};
\end{tikzpicture}

{\small Loading one page. Each arrow pair is one \omterm{def:g10:http-clients-servers:http}{request} and its \omterm{def:g10:http-clients-servers:http}{response}; time runs downwards. The browser cannot
ask for the stylesheet before it has read the page that names it.}
\end{omfigure}

\section{Forms and parameters}\label{sec:g10:http-clients-servers:forms}

\begin{definition}[Form, parameter]\label{def:g10:http-clients-servers:form}
A \emph{form}\index{form} is the part of a web page where the user types or chooses values and sends them; each value
is sent as a \emph{parameter}\index{parameter}, a pair \code{name=value}, to the address that the form names.
\end{definition}

\omcode{code/grade-10/08-http-clients-servers/listings/site_files/page.html}{1}{16}{The library page, with its
\omterm{def:g10:http-clients-servers:form}{form}.}\label{lst:g10:http-clients-servers:page}

\begin{omfigure}
\begin{tikzpicture}
  \omwindow{8}{4.2}{Brenmoor Library}
  \pic at (0,0) {window};
  \node[anchor=north west, font=\sffamily\bfseries\large] at (0.3,3.55) {Brenmoor Library};
  \fill[ocBlue] (0.75,2.25) circle (0.32); \node[font=\tiny\sffamily, text=white] at (0.75,2.25) {logo};
  \node[anchor=west, font=\sffamily\small] at (0.3,1.4) {Title:};
  \draw[black!60, fill=white] (1.3,1.18) rectangle (4.3,1.62);
  \node[draw=black!60, fill=black!8, rounded corners=1pt, font=\sffamily\small] at (5.0,1.4) {Search};
  \node[anchor=west, font=\sffamily\small] at (0.3,0.55) {Open every day.};
  \fill[ocBlue] (3.2,0.55) circle (0.22); \fill[ocBlue] (3.75,0.55) circle (0.22);
\end{tikzpicture}

{\small The page as a browser shows it. Typing \textit{forest} and pressing \textit{Search} makes the browser \omterm{def:g10:http-clients-servers:http}{request}
\code{/search?title=forest}: the \omterm{def:g10:http-clients-servers:form}{form}'s \code{method} is \code{get}.}
\end{omfigure}

\begin{definition}[Methods GET and POST]\label{def:g10:http-clients-servers:methods}
The \emph{method}\index{method (HTTP)} of a \omterm{def:g10:http-clients-servers:http}{request} says what the client wants. \emph{GET}\index{GET} asks for a
resource; its \omterm{def:g10:http-clients-servers:form}{parameters}, if any, travel in the \omterm{def:g10:http-clients-servers:url}{query string} of the address. \emph{POST}\index{POST} sends data to
the server, such as a \omterm{def:g10:http-clients-servers:form}{form} to record; its \omterm{def:g10:http-clients-servers:form}{parameters} travel in the body of the \omterm{def:g10:http-clients-servers:http}{request}.
\end{definition}

\omcode{code/grade-10/08-http-clients-servers/listings/library_site.py}{46}{57}{The site: a \omterm{def:g10:http-clients-servers:methods}{POST} \omterm{def:g10:http-clients-servers:http}{request}, whose
\omterm{def:g10:http-clients-servers:form}{parameters} are read from the body.}\label{lst:g10:http-clients-servers:site-post}

\omcode{code/grade-10/08-http-clients-servers/listings/form_get_post.py}{1}{12}{The same search, sent with \omterm{def:g10:http-clients-servers:methods}{GET} and
with \omterm{def:g10:http-clients-servers:methods}{POST}.}\label{lst:g10:http-clients-servers:form-get-post}

\omoutput{code/grade-10/08-http-clients-servers/outputs/form_get_post.txt}{With \omterm{def:g10:http-clients-servers:methods}{GET} the \omterm{def:g10:http-clients-servers:form}{parameter} is in the \omterm{def:g10:http-clients-servers:http}{request}
line, so it ends up in the server's log; with \omterm{def:g10:http-clients-servers:methods}{POST} it is in the body.}\label{out:g10:http-clients-servers:form-get-post}

% ledger: g10:http-9110-safe
\begin{method}[Choosing between GET and POST]\label{met:g10:http-clients-servers:get-post}
\begin{enumerate}
  \item A \omterm{def:g10:http-clients-servers:http}{request} that only reads (a search, a page) uses \omterm{def:g10:http-clients-servers:methods}{GET}: its address can be copied, bookmarked and sent again
    without harm. The rules of \omterm{def:g10:http-clients-servers:http}{HTTP} call \omterm{def:g10:http-clients-servers:methods}{GET} a \emph{safe} \omterm{def:g10:http-clients-servers:methods}{method}: it must not change anything on the server.
  \item A \omterm{def:g10:http-clients-servers:http}{request} that changes something (a loan, a message, an order) uses \omterm{def:g10:http-clients-servers:methods}{POST}, so that a crawler following links
    or a reloaded page does not repeat it unseen.
  \item Large data (a file, a long text) goes in a \omterm{def:g10:http-clients-servers:methods}{POST} body: addresses have a limited length.
  \item A password never goes in an address: addresses are logged and kept in the browser's history. It goes in a
    \omterm{def:g10:http-clients-servers:methods}{POST} body, over \omterm{def:g10:http-clients-servers:https}{HTTPS}.
\end{enumerate}
\end{method}

% ledger: g10:uri-3986
An address may only hold some \omterm{def:g10:text-encodings:ascii}{ASCII} characters, so the others are written by \emph{percent-encoding}: each byte of
their \omterm{def:g10:text-encodings:utf8}{UTF-8} encoding becomes \texttt{\%} followed by two hexadecimal digits. In a \omterm{def:g10:http-clients-servers:form}{form}'s \omterm{def:g10:http-clients-servers:url}{query string}, a space may
also be written \code{+}.

\omcode{code/grade-10/08-http-clients-servers/listings/percent.py}{1}{5}{Percent-encoding and
back.}\label{lst:g10:http-clients-servers:percent}

\omoutput{code/grade-10/08-http-clients-servers/outputs/percent.txt}{\textit{é} is the two bytes \code{C3 A9}; the
space is \code{20} and \code{&} is \code{26}.}\label{out:g10:http-clients-servers:percent}

\section{Cookies and sessions}\label{sec:g10:http-clients-servers:cookies}

Each \omterm{def:g10:http-clients-servers:http}{HTTP} \omterm{def:g10:http-clients-servers:http}{request} stands alone: the server does not remember the previous one. To recognise a visitor, it hands the
browser a small piece of text that the browser sends back.

% ledger: g10:cookie-6265
\begin{definition}[Cookie, session]\label{def:g10:http-clients-servers:cookie}
A \emph{cookie}\index{cookie} is a pair \code{name=value} that a server sends in a \code{Set-Cookie} \omterm{def:g10:http-clients-servers:http-header}{header} and that
the browser stores and sends back, in a \code{Cookie} \omterm{def:g10:http-clients-servers:http-header}{header}, with every later \omterm{def:g10:http-clients-servers:http}{request} to the same server. A
\emph{session}\index{session} is the series of \omterm{def:g10:http-clients-servers:http}{requests} of one visitor that a server links together, usually through
a cookie holding a random \emph{session identifier}, while the session's data stays on the server.
\end{definition}

\omcode{code/grade-10/08-http-clients-servers/listings/library_site.py}{59}{63}{The site: a visit counter kept in a
\omterm{def:g10:http-clients-servers:cookie}{cookie}.}\label{lst:g10:http-clients-servers:site-count}

\omcode{code/grade-10/08-http-clients-servers/listings/cookie_counter.py}{1}{13}{A client with a \omterm{def:g10:http-clients-servers:cookie}{cookie} jar, like a
browser.}\label{lst:g10:http-clients-servers:cookie-counter}

\omoutput{code/grade-10/08-http-clients-servers/outputs/cookie_counter.txt}{Three visits: each \omterm{def:g10:http-clients-servers:http}{response} sets the
\omterm{def:g10:http-clients-servers:cookie}{cookie} that the next \omterm{def:g10:http-clients-servers:http}{request} sends back.}\label{out:g10:http-clients-servers:cookie-counter}

The counter lives in the browser: a visitor could edit the \omterm{def:g10:http-clients-servers:cookie}{cookie} and claim a thousand visits. Anything the server
must trust, such as who is logged in, stays on the server, and the \omterm{def:g10:http-clients-servers:cookie}{cookie} holds only an identifier too long and
random to guess (\cref{exo:g10:http-clients-servers:13}).

\begin{safety}
\omterm{def:g10:http-clients-servers:cookie}{Cookies} are also how a site, or an advertiser present on many sites, recognises the same browser from visit to
visit. Browsers let you see and delete the \omterm{def:g10:http-clients-servers:cookie}{cookies} of each site, block the \omterm{def:g10:http-clients-servers:cookie}{cookies} of sites other than the one you
are visiting, and delete everything when the window is closed. A private window starts with no \omterm{def:g10:http-clients-servers:cookie}{cookies} and forgets
them at the end.
\end{safety}

\section{What runs where}\label{sec:g10:http-clients-servers:where}

The server's code ran on the server: it chose the answer, and the browser received only the result. A page can also
carry code that runs in the browser, written in the language \emph{JavaScript}. Such code reacts to
\emph{events}, a click on a button, a key typed in a box, the end of loading, by running a function called a
\emph{handler}: it can check a \omterm{def:g10:http-clients-servers:form}{form} before it is sent, or change part of the page without loading a new one.
Whatever runs in the browser can be read and changed by the user, so a server checks every value it receives again.

\begin{definition}[HTTPS]\label{def:g10:http-clients-servers:https}
\emph{HTTPS}\index{HTTPS} is \omterm{def:g10:http-clients-servers:http}{HTTP} sent over an encrypted connection: the same \omterm{def:g10:http-clients-servers:http}{requests} and \omterm{def:g10:http-clients-servers:http}{responses}, but anyone who
watches the packets on the way sees only unreadable bytes, and the browser checks that the server really is the one
named in the address. Grade 12 explains how.
\end{definition}

\begin{history}[The web]
% ledger: g10:web-1989, g10:web-1990
The web was proposed in March 1989 by Tim Berners-Lee, then working at a European physics laboratory, as a way to
share documents linked to one another. By the end of 1990 the first web server and the first browser ran there, and
the address of the first site was \code{info.cern.ch}. \omterm{def:g10:http-clients-servers:http}{HTTP}, HTML and \omterm{def:g10:http-clients-servers:url}{URLs} were the three pieces of the proposal,
and they still are.
\end{history}

\section{Exercises}\label{sec:g10:http-clients-servers:exercises}

\begin{exercise}[$\star$]\label{exo:g10:http-clients-servers:1}
Give the \omterm{def:g10:http-clients-servers:url}{scheme}, \omterm{def:g10:http-clients-servers:url}{host}, port, path, \omterm{def:g10:http-clients-servers:url}{query string} and \omterm{def:g10:http-clients-servers:url}{fragment} of
\texttt{https://maps.brenmoor.example/route?from=3\&to=17\#map}. Which port does the browser use?
\end{exercise}

\begin{exercise}[$\star$]\label{exo:g10:http-clients-servers:2}
Give the class of each \omterm{def:g10:http-clients-servers:status}{status code}, and say whose fault each error is: \code{201}, \code{301}, \code{403},
\code{500}, \code{503}.
\end{exercise}

\begin{exercise}[$\star$]\label{exo:g10:http-clients-servers:3}
The page's \omterm{def:g10:http-clients-servers:form}{form} has \code{method="get"}. A user types \textit{the lost code} and presses \textit{Search}. Write the
\omterm{def:g10:http-clients-servers:http}{request} line the browser sends.
\end{exercise}

\begin{exercise}[$\star$]\label{exo:g10:http-clients-servers:4}
A browser visits \code{/count} twice, then deletes the site's \omterm{def:g10:http-clients-servers:cookie}{cookies}, then visits it once more. What does each
\omterm{def:g10:http-clients-servers:http}{response} say, and what is in the jar at the end?
\end{exercise}

\begin{exercise}[$\star$]\label{exo:g10:http-clients-servers:5}
Write \code{query_params(url)}, which returns the dictionary of the \omterm{def:g10:http-clients-servers:form}{parameters} of an address's \omterm{def:g10:http-clients-servers:url}{query string}, without
\code{urllib}. Leave the values percent-encoded.
\end{exercise}

\begin{exercise}[$\star\star$]\label{exo:g10:http-clients-servers:6}
A student writes a \omterm{def:g10:http-clients-servers:http}{response} by hand as a status line, a \code{Content-Type} \omterm{def:g10:http-clients-servers:http-header}{header}, then the body. The client waits
and never shows the page. Find the two things missing, and say what each one tells the client.
\end{exercise}

\begin{exercise}[$\star\star$]\label{exo:g10:http-clients-servers:7}
In the site, the last branch of \code{do_GET} answers \code{404}. What would happen to a \omterm{def:g10:http-clients-servers:http}{request} for an unknown path
if that branch were missing? Why is an answer always needed?
\end{exercise}

\begin{exercise}[$\star\star$]\label{exo:g10:http-clients-servers:8}
Count the bytes of the \omterm{def:g10:http-clients-servers:http}{request} of the raw exchange listing, the \texttt{\textbackslash r\textbackslash n} included. How many of them are
\omterm{def:g10:http-clients-servers:http-header}{headers} and empty line, and how many \omterm{def:g10:http-clients-servers:http}{request} line?
\end{exercise}

\begin{exercise}[$\star\star$]\label{exo:g10:http-clients-servers:9}
A login \omterm{def:g10:http-clients-servers:form}{form} uses \code{method="get"} over \code{http}. The address after logging in is
\code{/login?user=noor&password=tulip7}. List three places where the password can now be read, and correct the
\omterm{def:g10:http-clients-servers:form}{form}.
\end{exercise}

\begin{exercise}[$\star\star$]\label{exo:g10:http-clients-servers:10}
Percent-encode \textit{café \& co} by hand, as in a path (space as \texttt{\%20}). Then explain why a raw \code{&}
cannot be left inside a \omterm{def:g10:http-clients-servers:url}{query string}'s value.
\end{exercise}

\begin{exercise}[$\star\star$]\label{exo:g10:http-clients-servers:11}
Add a route \code{/sum?a=2&b=40} to the site that answers the sum, \code{42}, and answers \code{400} (bad \omterm{def:g10:http-clients-servers:http}{request})
when a \omterm{def:g10:http-clients-servers:form}{parameter} is missing or not a whole number.
\end{exercise}

\begin{exercise}[$\star\star\star$]\label{exo:g10:http-clients-servers:12}
A page names $1$ stylesheet and $3$ images. Each new connection costs one round trip before the first \omterm{def:g10:http-clients-servers:http}{request}, and
each \omterm{def:g10:http-clients-servers:http}{request} and its \omterm{def:g10:http-clients-servers:http}{response} cost one more. How many round trips does loading the page take with a new connection
for every \omterm{def:g10:http-clients-servers:http}{request}? With a single connection reused? Why can the stylesheet never be requested in the same round trip
as the page?
\end{exercise}

\begin{exercise}[$\star\star\star$]\label{exo:g10:http-clients-servers:13}
Write a route \code{/basket?add=bread} that keeps a shopping basket per visitor on the server: a dictionary from a
\omterm{def:g10:http-clients-servers:cookie}{session} identifier to the list of items, the identifier made with \code{secrets.token_hex(8)} and sent in a \omterm{def:g10:http-clients-servers:cookie}{cookie}.
Why is it safer than keeping the basket itself in the \omterm{def:g10:http-clients-servers:cookie}{cookie}?
\end{exercise}

\begin{exercise}[$\star\star\star$]\label{exo:g10:http-clients-servers:14}
A route \code{/a} redirects to \code{/b}, and \code{/b} redirects to \code{/a}. What does a client do? Try it with
\code{urllib}, and explain how a client can detect such a loop.
\end{exercise}

\begin{exercise}[$\star\star\star$]\label{exo:g10:http-clients-servers:15}
A site deletes a book from its catalogue when \code{GET /delete?book=12} is requested, from a link on an
administrator's page. Describe how a search engine's crawler, which follows every link it finds, could empty the
catalogue, and propose a correct design.
\end{exercise}

\section{Problem: What One Page Costs}\label{sec:g10:http-clients-servers:problem}

\begin{problem}[{Weekend problem --- how long the library page takes to load}]\label{pb:g10:http-clients-servers:1}
The program \code{exercises/page_cost.py} starts the library site, \omterm{def:g10:http-clients-servers:http}{requests} \code{page.html} as a browser would,
reads it, \omterm{def:g10:http-clients-servers:http}{requests} every resource it names, and prints the size of every body. Take a latency of $50$~ms: a round
trip, a message there and its answer back, takes $50$~ms. Opening a connection costs one round trip; each \omterm{def:g10:http-clients-servers:http}{request}
and its \omterm{def:g10:http-clients-servers:http}{response} cost one more.

\medskip\noindent\textbf{Part I --- One page and its resources.}
\begin{enumerate}
  \item Read \code{page.html}: which resources does it name, and with which tags?
  \item How many \omterm{def:g10:http-clients-servers:http}{requests} does loading the page take in all?
  \item Which \omterm{def:g10:http-clients-servers:http}{request} must come first, and why?
  \item Which \omterm{def:g10:http-clients-servers:status}{status code} does each of them receive?
\end{enumerate}

\medskip\noindent\textbf{Part II --- Requests and bytes.}
\begin{enumerate}[resume]
  \item Run the program. Give the size of each body and the total.
  \item The \omterm{def:g10:http-clients-servers:http}{response} \omterm{def:g10:http-clients-servers:http-header}{headers} of the site take about $130$ bytes per \omterm{def:g10:http-clients-servers:http}{response}. Estimate the total number of bytes
    received, \omterm{def:g10:http-clients-servers:http-header}{headers} included. Which part of it is \omterm{def:g10:http-clients-servers:http-header}{headers}?
  \item The \omterm{def:g10:http-clients-servers:http}{requests} a browser sends are about $400$ bytes each, with many \omterm{def:g10:http-clients-servers:http-header}{headers}. Compare the bytes sent and the
    bytes received.
  \item At $10$ megabits per second, how long does it take to receive the bytes of question 6? Compare with one round
    trip.
\end{enumerate}

\medskip\noindent\textbf{Part III --- Round trips.}
\begin{enumerate}[resume]
  \item With a new connection for each \omterm{def:g10:http-clients-servers:http}{request}, how many round trips does the page take, and how many milliseconds?
  \item With one connection reused for every \omterm{def:g10:http-clients-servers:http}{request}?
  \item A browser opens up to six connections at once to one server. With two connections opened together after the
    HTML has arrived, how many round trips does the page take?
  \item Which matters more for this page: the number of bytes or the number of round trips?
  \item With one reused connection, what latency would make the page take a whole second?
\end{enumerate}

\medskip\noindent\textbf{Part IV --- Caching.}
\begin{enumerate}[resume]
  \item On a second visit, the browser still holds the four resources. If it asks for each one again with a \omterm{def:g10:http-clients-servers:http-header}{header}
    saying which copy it has, the server can answer \code{304 Not Modified} with no body. How many round trips does
    that visit take with one connection?
  \item If the server says, in a \omterm{def:g10:http-clients-servers:http-header}{header} of the first visit, that the resources stay valid for a day, the browser
    does not ask at all. How many round trips then?
  \item Why can the HTML page itself rarely be kept that way?
  \item A reader opens the page ten times a day. With one reused connection, how many round trips per day does the
    day-long validity save, compared with asking again each time as in question 14?
  \item State the result: the \omterm{def:g10:http-clients-servers:http}{requests}, the bytes and the round trips of the page, and its load time with a new
    connection per \omterm{def:g10:http-clients-servers:http}{request} and with one reused connection.
\end{enumerate}
\end{problem}
